% chapter06.tex — Continuous Functions
% Book pp. 115–121. Scan: sources/chapters/ch06-continuous-functions.pdf
% (slice of original pp. 136–142). PDF pages 4 and 5 are swapped
% (book 119 then 118); this file is in printed order.

\spivakchapter{6}{Continuous Functions}

If $f$ is an arbitrary function, it is not necessarily true that
\[
\lim_{x \to a} f(x) = f(a).
\]
In fact, there are many ways this can fail to be true. For example, $f$
might not even be defined at $a$, in which case the equation makes no
sense (Figure~1).

\begin{center}
\figtex{assets/geometric/ch06-fig01}\\[0.6em]
\figlabel{1}
\end{center}

Again, $\lim_{x \to a} f(x)$ might not exist (Figure~2). Finally, as
illustrated in Figure~3, even if $f$ is defined at $a$ and
$\lim_{x \to a} f(x)$ exists, the limit might not equal $f(a)$.

\begin{center}
\figtex{assets/geometric/ch06-fig02}\\[0.6em]
\figlabel{2}
\end{center}

\begin{center}
\figtex{assets/geometric/ch06-fig03}\\[0.6em]
\figlabel{3}
\end{center}

We would like to regard all behavior of this type as abnormal and honor,
with some complimentary designation, functions which do not exhibit such
peculiarities. The term which has been adopted is ``continuous.''
Intuitively, a function $f$ is continuous if the graph contains no
breaks, jumps, or wild oscillations. Although this description will
usually enable you to decide whether a function is continuous simply by
looking at its graph (a skill well worth cultivating) it is easy to be
fooled, and the precise definition is \emph{very} important.

\begin{spivakdefn}
The function $f$ is \textbf{continuous at $a$} if
\[
\lim_{x \to a} f(x) = f(a).
\]
\end{spivakdefn}

We will have no difficulty finding many examples of functions which are,
or are not, continuous at some number $a$---every example involving
limits provides an example about continuity, and Chapter~5 certainly
provides enough of these.

The function $f(x) = \sin 1/x$ is not continuous at $0$, because it is
not even defined at $0$, and the same is true of the function
$g(x) = x \sin 1/x$. On the other hand, if we are willing to extend the
second of these functions, that is, if we wish to define a new function
$G$ by
\[
G(x) =
\begin{cases}
x \sin 1/x, & x \neq 0 \\
a,          & x = 0,
\end{cases}
\]
then the choice of $a = G(0)$ can be made in such a way that $G$ will be
continuous at $0$---to do this we can (if fact, we must) define
$G(0) = 0$ (Figure~4). This sort of extension is not possible for $f$;
if we define
\[
F(x) =
\begin{cases}
\sin 1/x, & x \neq 0 \\
a,        & x = 0,
\end{cases}
\]
then $F$ will not be continuous at $0$, no matter what $a$ is, because
$\lim_{x \to 0} f(x)$ does not exist.

\begin{center}
\figtex{assets/plots/ch06-fig04}\\[0.6em]
\figlabel{4}
\end{center}

The function
\[
f(x) =
\begin{cases}
x, & x \text{ rational} \\
0, & x \text{ irrational}
\end{cases}
\]
is not continuous at $a$, if $a \neq 0$, since $\lim_{x \to a} f(x)$
does not exist. However, $\lim_{x \to 0} f(x) = 0 = f(0)$, so $f$ is
continuous at precisely one point, $0$.

The functions $f(x) = c$, $g(x) = x$, and $h(x) = x^2$ are continuous at
all numbers $a$, since
\begin{align*}
\lim_{x \to a} f(x) &= \lim_{x \to a} c = c = f(a), \\
\lim_{x \to a} g(x) &= \lim_{x \to a} x = a = g(a), \\
\lim_{x \to a} h(x) &= \lim_{x \to a} x^2 = a^2 = h(a).
\end{align*}

Finally, consider the function
\[
f(x) =
\begin{cases}
0,    & x \text{ irrational} \\
1/q,  & x = p/q \text{ in lowest terms.}
\end{cases}
\]
In Chapter~5 we showed that $\lim_{x \to a} f(x) = 0$ for all $a$
(actually, only for $0 < a < 1$, but you can easily see that this is
true for all $a$). Since $0 = f(a)$ only when $a$ is irrational, this
function is continuous at $a$ if $a$ is irrational, but not if $a$ is
rational.

It is even easier to give examples of continuity if we prove two simple
theorems.

\begin{theorem}
If $f$ and $g$ are continuous at $a$, then
\begin{steps}
\item $f + g$ is continuous at $a$,
\item $f \cdot g$ is continuous at $a$.
\end{steps}
Moreover, if $g(a) \neq 0$, then
\begin{steps}[start=3]
\item $1/g$ is continuous at $a$.
\end{steps}
\end{theorem}

\begin{proof}
Since $f$ and $g$ are continuous at $a$,
\[
\lim_{x \to a} f(x) = f(a)
\quad\text{and}\quad
\lim_{x \to a} g(x) = g(a).
\]
By Theorem~2(1) of Chapter~5 this implies that
\[
\lim_{x \to a} (f + g)(x) = f(a) + g(a) = (f + g)(a),
\]
which is just the assertion that $f + g$ is continuous at $a$. The
proofs of parts~(2) and~(3) are left to you.
\end{proof}

Starting with the functions $f(x) = c$ and $f(x) = x$, which are
continuous at $a$, for every $a$, we can use Theorem~1 to conclude that
a function
\[
f(x) = \frac{b_n x^n + b_{n-1} x^{n-1} + \cdots + b_0}
{c_m x^m + c_{m-1} x^{m-1} + \cdots + c_0}
\]
is continuous at every point in its domain. But it is harder to get
much further than that. When we discuss the sine function in detail it
will be easy to prove that $\sin$ is continuous at $a$ for all $a$; let
us assume this fact meanwhile. A function like
\[
f(x) = \frac{\sin^2 x + x^2 + x^4 \sin x}
{\sin^{27} x + 4x^2 \sin^2 x}
\]
can now be proved continuous at every point in its domain. But we are
still unable to prove the continuity of a function like
$f(x) = \sin(x^2)$; we obviously need a theorem about the composition of
continuous functions. Before stating this theorem, the following point
about the definition of continuity is worth noting. If we translate the
equation $\lim_{x \to a} f(x) = f(a)$ according to the definition of
limits, we obtain
\begin{quote}
for every $\varepsilon > 0$ there is $\delta > 0$ such that, for all
$x$,\\
if $0 < \abs{x - a} < \delta$, then $\abs{f(x) - f(a)} < \varepsilon$.
\end{quote}
But in this case, where the limit is $f(a)$, the phrase
\[
0 < \abs{x - a} < \delta
\]
may be changed to the simpler condition
\[
\abs{x - a} < \delta,
\]
since if $x = a$ it is certainly true that
$\abs{f(x) - f(a)} < \varepsilon$.

\begin{theorem}
If $g$ is continuous at $a$, and $f$ is continuous at $g(a)$, then
$f \circ g$ is continuous at $a$. (Notice that $f$ is required to be
continuous at $g(a)$, not at $a$.)
\end{theorem}

\begin{proof}
Let $\varepsilon > 0$. We wish to find a $\delta > 0$ such that for all
$x$,
\begin{align*}
&\text{if } \abs{x - a} < \delta, \text{ then }
  \abs{(f \circ g)(x) - (f \circ g)(a)} < \varepsilon, \\
&\qquad\text{i.e., } \abs{f(g(x)) - f(g(a))} < \varepsilon.
\end{align*}
We first use continuity of $f$ to estimate how close $g(x)$ must be to
$g(a)$ in order for this inequality to hold. Since $f$ is continuous at
$g(a)$, there is a $\delta' > 0$ such that for all $y$,
\[
(1)\qquad
\text{if } \abs{y - g(a)} < \delta',\ \text{then }
\abs{f(y) - f(g(a))} < \varepsilon.
\]
In particular, this means that
\[
(2)\qquad
\text{if } \abs{g(x) - g(a)} < \delta',\ \text{then }
\abs{f(g(x)) - f(g(a))} < \varepsilon.
\]
We now use continuity of $g$ to estimate how close $x$ must be to $a$
in order for the inequality $\abs{g(x) - g(a)} < \delta'$ to hold. The
number $\delta'$ is a positive number just like any other positive
number; we can therefore take $\delta'$ as the $\varepsilon$~(!) in the
definition of continuity of $g$ at $a$. We conclude that there is a
$\delta > 0$ such that, for all $x$,
\[
(3)\qquad
\text{if } \abs{x - a} < \delta,\ \text{then }
\abs{g(x) - g(a)} < \delta'.
\]
Combining~(2) and~(3) we see that for all $x$,
\[
\text{if } \abs{x - a} < \delta,\ \text{then }
\abs{f(g(x)) - f(g(a))} < \varepsilon.
\]
\end{proof}

We can now reconsider the function
\[
f(x) =
\begin{cases}
x \sin 1/x, & x \neq 0 \\
0,          & x = 0.
\end{cases}
\]
We have already noted that $f$ is continuous at $0$. A few applications
of Theorems~1 and~2, together with the continuity of $\sin$, show that
$f$ is also continuous at $a$, for $a \neq 0$. Functions like
$f(x) = \sin\bigl(x^2 + \sin(x + \sin^2(x^3))\bigr)$ should be equally
easy for you to analyze.

The few theorems of this chapter have all been related to continuity of
functions at a single point, but the concept of continuity doesn't
begin to be really interesting until we focus our attention on functions
which are continuous at all points of some interval. If $f$ is
continuous at $x$ for all $x$ in $(a,b)$, then $f$ is called
\textbf{continuous on $(a,b)$}; as a ``special case'', $f$ is
\textbf{continuous on $\mathbf{R} = (-\infty,\infty)$} [see page~57] if
it is continuous at $x$ for all $x$ in $\mathbf{R}$. Continuity on a
closed interval must be defined a little differently; a function $f$ is
called \textbf{continuous on $[a,b]$} if
\begin{steps}
\item $f$ is continuous at $x$ for all $x$ in $(a,b)$,
\item $\lim_{x \to a^+} f(x) = f(a)$ and $\lim_{x \to b^-} f(x) = f(b)$.
\end{steps}
(We also often simply say that a function is continuous if it is
continuous at $x$ for all $x$ in its domain.)

Functions which are continuous on an interval are usually regarded as
especially well behaved; indeed continuity might be specified as the
first condition which a ``reasonable'' function ought to satisfy. A
continuous function is sometimes described, intuitively, as one whose
graph can be drawn without lifting your pencil from the paper.
Consideration of the function
\[
f(x) =
\begin{cases}
x \sin 1/x, & x \neq 0 \\
0,          & x = 0
\end{cases}
\]
shows that this description is a little too optimistic, but it is
nevertheless true that there are many important results involving
functions which are continuous on an interval. There theorems are
generally much harder than the ones in this chapter, but there is a
simple theorem which forms a bridge between the two kinds of results.
The hypothesis of this theorem requires continuity at only a single
point, but the conclusion describes the behavior of the function on some
interval containing the point. Although this theorem is really a lemma
for later arguments, it is included here as a preview of things to come.

\begin{theorem}
Suppose $f$ is continuous at $a$, and $f(a) > 0$. Then $f(x) > 0$ for
all $x$ in some interval containing $a$; more precisely, there is a
number $\delta > 0$ such that $f(x) > 0$ for all $x$ satisfying
$\abs{x - a} < \delta$. Similarly, if $f(a) < 0$, then there is a
number $\delta > 0$ such that $f(x) < 0$ for all $x$ satisfying
$\abs{x - a} < \delta$.
\end{theorem}

\begin{proof}
Consider the case $f(a) > 0$. Since $f$ is continuous at $a$, for every
$\varepsilon > 0$ there is a $\delta > 0$ such that, for all $x$,
\begin{align*}
&\text{if } \abs{x - a} < \delta, \text{ then }
  \abs{f(x) - f(a)} < \varepsilon, \\
&\qquad\text{i.e., } {-\varepsilon} < f(x) - f(a) < \varepsilon.
\end{align*}
In particular, this must hold for $\varepsilon = \frac12 f(a)$, since
$\frac12 f(a) > 0$ (Figure~5). Thus there is $\delta > 0$ so that for
all $x$,
\[
\text{if } \abs{x - a} < \delta, \text{ then }
{-\tfrac12 f(a)} < f(x) - f(a) < \tfrac12 f(a),
\]
and this implies that $f(x) > \frac12 f(a) > 0$. (We could even have
picked $\varepsilon$ to be $f(a)$ itself, leading to a proof that is
more elegant, but more confusing to picture.)

\begin{center}
\figtex{assets/geometric/ch06-fig05}\\[0.6em]
\figlabel{5}
\end{center}

A similar proof can be given in the case $f(a) < 0$; take
$\varepsilon = -\frac12 f(a)$. Or one can apply the first case to the
function $-f$.
\end{proof}

\subsection*{PROBLEMS}

\pnum{1.}
For which of the following functions $f$ is there a continuous function
$F$ with domain $\mathbf{R}$ such that $F(x) = f(x)$ for all $x$ in the
domain of $f$?
\begin{romans}
\item $f(x) = \dfrac{x^2 - 4}{x - 2}$.
\item $f(x) = \dfrac{\abs{x}}{x}$.
\item $f(x) = 0$, $x$ irrational.
\item $f(x) = 1/q$, $x = p/q$ rational in lowest terms.
\end{romans}

\pnum{2.}
At which points are the functions of Problems~4-17 and~4-19 continuous?

\pnum{3.}
\begin{letters}
\item Suppose that $f$ is a function satisfying $\abs{f(x)} \leq \abs{x}$
  for all $x$. Show that $f$ is continuous at $0$. (Notice that $f(0)$
  must equal $0$.)
\item Give an example of such a function $f$ which is not continuous at
  any $a \neq 0$.
\item Suppose that $g$ is continuous at $0$ and $g(0) = 0$, and
  $\abs{f(x)} \leq \abs{g(x)}$. Prove that $f$ is continuous at $0$.
\end{letters}

\pnum{4.}
Give an example of a function $f$ such that $f$ is continuous nowhere,
but $\abs{f}$ is continuous everywhere.

\pnum{5.}
For each number $a$, find a function which is continuous at $a$, but not
at any other points.

\pnum{6.}
\begin{letters}
\item Find a function $f$ which is discontinuous at
  $1, \frac12, \frac13, \frac14, \dots$ but continuous at all other
  points.
\item Find a function $f$ which is discontinuous at
  $1, \frac12, \frac13, \frac14, \dots$, and at $0$, but continuous at
  all other points.
\end{letters}

\pnum{7.}
Suppose that $f$ satisfies $f(x + y) = f(x) + f(y)$, and that $f$ is
continuous at $0$. Prove that $f$ is continuous at $a$ for all $a$.

\pnum{8.}
Suppose that $f$ is continuous at $a$ and $f(a) = 0$. Prove that if
$\alpha \neq 0$, then $f + \alpha$ is nonzero in some open interval
containing $a$.

\pnum{9.}
\begin{letters}
\item Suppose $f$ is defined at $a$ but is not continuous at $a$. Prove
  that for some number $\varepsilon > 0$ there are numbers $x$
  arbitrarily close to $a$ with $\abs{f(x) - f(a)} > \varepsilon$.
  Illustrate graphically.
\item Conclude that for some number $\varepsilon > 0$ \emph{either}
  there are numbers $x$ arbitrarily close to $a$ with
  $f(x) < f(a) - \varepsilon$ \emph{or} there are numbers $x$
  arbitrarily close to $a$ with $f(x) > f(a) + \varepsilon$.
\end{letters}

\pnum{10.}
\begin{letters}
\item Prove that if $f$ is continuous at $a$, then so is $\abs{f}$.
\item Prove that every function $f$ continuous on $\mathbf{R}$ can be
  written $f = E + O$, where $E$ is even and continuous and $O$ is odd
  and continuous.
\item Prove that if $f$ and $g$ are continuous, then so are
  $\max(f,g)$ and $\min(f,g)$.
\item Prove that every continuous $f$ can be written $f = g - h$, where
  $g$ and $h$ are nonnegative and continuous.
\end{letters}

\pnum{11.}
Prove Theorem~1(3) by using Theorem~2 and continuity of the function
$f(x) = 1/x$.

\pnum{*12.}
\begin{letters}
\item Prove that if $f$ is continuous at $l$ and
  $\lim_{x \to a} g(x) = l$, then $\lim_{x \to a} f(g(x)) = f(l)$. (You
  can go right back to the definitions, but it is easier to consider the
  function $G$ with $G(x) = g(x)$ for $x \neq a$, and $G(a) = l$.)
\item Show that if continuity of $f$ at $l$ is not assumed, then it is
  not generally true that
  \[
  \lim_{x \to a} f(g(x)) = f\bigl(\lim_{x \to a} g(x)\bigr).
  \]
  \textit{Hint:} Try $f(x) = 0$ for $x \neq l$, and $f(l) = 1$.
\end{letters}

\pnum{13.}
\begin{letters}
\item Prove that if $f$ is continuous on $[a,b]$, then there is a
  function $g$ which is continuous on $\mathbf{R}$, and which satisfies
  $g(x) = f(x)$ for all $x$ in $[a,b]$. \textit{Hint:} Since you
  obviously have a great deal of choice, try making $g$ constant on
  $(-\infty, a]$ and $[b, \infty)$.
\item Give an example to show that this assertion is false if $[a,b]$
  is replaced by $(a,b)$.
\end{letters}

\pnum{14.}
\begin{letters}
\item Suppose that $g$ and $h$ are continuous at $a$, and that
  $g(a) = h(a)$. Define $f(x)$ to be $g(x)$ if $x \geq a$ and $h(x)$ if
  $x \leq a$. Prove that $f$ is continuous at $a$.
\item Suppose $g$ is continuous on $[a,b]$ and $h$ is continuous on
  $[b,c]$ and $g(b) = h(b)$. Let $f(x)$ be $g(x)$ for $x$ in $[a,b]$ and
  $h(x)$ for $x$ in $[b,c]$. Show that $f$ is continuous on $[a,c]$.
  (Thus, continuous functions can be ``pasted together''.)
\end{letters}

\pnum{15.}
Prove that if $f$ is continuous at $a$, then for any $\varepsilon > 0$
there is a $\delta > 0$ so that whenever $\abs{x - a} < \delta$ and
$\abs{y - a} < \delta$, we have $\abs{f(x) - f(y)} < \varepsilon$.

\pnum{16.}
\begin{letters}
\item Prove the following version of Theorem~3 for ``right-hand
  continuity'': Suppose that
  \[
  \lim_{x \to a^+} f(x) = f(a),
  \]
  and $f(a) > 0$. Then there is a number $\delta > 0$ such that
  $f(x) > 0$ for all $x$ satisfying $0 \leq x - a < \delta$. Similarly,
  if $f(a) < 0$, then there is a number $\delta > 0$ such that
  $f(x) < 0$ for all $x$ satisfying $0 \leq x - a < \delta$.
\item Prove a version of Theorem~3 when $\lim_{x \to b^-} f(x) = f(b)$.
\end{letters}

\pnum{17.}
If $\lim_{x \to a} f(x)$ exists, but is $\neq f(a)$, then $f$ is said to
have a \textbf{removable discontinuity} at $a$.
\begin{letters}
\item If $f(x) = \sin 1/x$ for $x \neq 0$ and $f(0) = 1$, does $f$ have
  a removable discontinuity at $0$? What if $f(x) = x \sin 1/x$ for
  $x \neq 0$, and $f(0) = 1$?
\item Suppose $f$ has a removable discontinuity at $a$. Let
  $g(x) = f(x)$ for $x \neq a$, and let $g(a) = \lim_{x \to a} f(x)$.
  Prove that $g$ is continuous at $a$. (Don't work very hard; this is
  quite easy.)
\item Let $f(x) = 0$ if $x$ is irrational, and let $f(p/q) = 1/q$ if
  $p/q$ is in lowest terms. What is the function $g$ defined by
  $g(x) = \lim_{y \to x} f(y)$?
\item[*(d)] Let $f$ be a function with the property that every point of
  discontinuity is a removable discontinuity. This means that
  $\lim_{y \to x} f(y)$ exists for all $x$, but $f$ may be discontinuous
  at some (even infinitely many) numbers $x$. Define
  $g(x) = \lim_{y \to x} f(y)$. Prove that $g$ is continuous. (This is
  not quite so easy as part~(b).)
\item[**(e)] Is there a function $f$ which is discontinuous at every
  point, and which has only removable discontinuities? (It is worth
  thinking about this problem now, but mainly as a test of intuition;
  even if you suspect the correct answer, you will almost certainly be
  unable to prove it at the present time. See Problem~22-33.)
\end{letters}
